\documentclass[twocolumn]{iopjournal}

\usepackage{amsmath,amssymb,amsfonts}
\usepackage{amsthm}

\fancyhfoffset[L]{0mm}
\setlength{\marginparwidth}{0pt}
\setlength{\marginparsep}{0pt}

% ==========================================
% ZENODO PREPRINT HEADERS (ACTIVE)
% ==========================================
\renewcommand{\articletype}[1]{%
  {\vspace*{-8mm}\noindent \Large \sf Preprint}\par
  \vspace*{6mm}\noindent{\scriptsize \sf{\bfseries \MakeUppercase{#1}}}}

\fancyhead[L]{}
\fancyhead[R]{}

% Theorem environments
\newtheorem{lemma}{Lemma}
\newtheorem{theorem}[lemma]{Theorem}
\newtheorem{corollary}[lemma]{Corollary}
\theoremstyle{definition}
\newtheorem{definition}[lemma]{Definition}
\newtheorem{remark}[lemma]{Remark}

\begin{document}

\articletype{Supplementary Material}

\title{The Energy Functional of the Gravitational Constitutive Law}

\author{Stephen Nelson}

\affil{Field Dynamics Research}

\ead{stephen@fielddynamics.org}

\begin{abstract}
The AQUAL framework of Bekenstein and Milgrom requires a kinetic
functional $F(y)$, with $y = |\nabla\Phi|^2/a_0^2$, whose derivative
with respect to $y$ is the interpolating function $\mu$.  For the
simple interpolating function $\mu(x) = x/(x+1)$, widely adopted in
MOND rotation-curve fits, $F$ has historically been carried
implicitly through its derivative.  We give the closed form
$F(y) = y - 2\sqrt{y} + 2\ln(1 + \sqrt{y})$, derived by elementary
integration and verified independently by differentiation.  The
result has three terms with distinct asymptotic roles: the
polynomial $y$ dominates at high acceleration and recovers
Newtonian behaviour, the logarithm dominates at low acceleration
and produces deep-MOND scaling, and the square root provides smooth
crossover with no piecewise patching.  The closed form is suitable
for direct use in action-level calculations where implicit
formulations of $F$ are not workable.
\end{abstract}

% ==================================================================

\section*{The AQUAL identity}

The AQUAL Lagrangian density of Bekenstein and
Milgrom~\cite{BekensteinMilgrom1984} is
\begin{equation}
  \mathcal{L}_{\text{AQUAL}}
  = -\frac{a_0^2}{8\pi G}\, F(y) + \rho\,\Phi,
  \qquad
  y = \frac{|\nabla\Phi|^2}{a_0^2},
  \label{eq:lagrangian}
\end{equation}
where $a_0$ is the MOND acceleration scale of
Milgrom~\cite{Milgrom1983}.  Variation with respect to $\Phi$ yields
the modified Poisson equation
\begin{equation}
  \nabla \cdot \left[\mu\!\left(\frac{|\nabla\Phi|}{a_0}\right)
  \nabla\Phi \right] = 4\pi G\rho,
  \label{eq:modified_poisson}
\end{equation}
in which the interpolating function $\mu$ is fixed by the AQUAL
identity
\begin{equation}
  \mu(\sqrt{y}) = \frac{dF}{dy}.
  \label{eq:aqual_identity}
\end{equation}

\noindent
The framework determines $\mu$ from $F$ by differentiation and $F$
from $\mu$ by integration, but the choice of either is left to the
user.  Empirically, the simple interpolating
function~\cite{FamaeyMcGaugh2012}
\begin{equation}
  \mu(x) = \frac{x}{x+1}
  \label{eq:simple_mu}
\end{equation}
provides good fits to galaxy rotation curves across a wide range of
masses.  No closed form for the corresponding $F(y)$ appears in
the published literature; this note supplies it.

% ==================================================================

\section*{Main result}

\begin{theorem}[Closed-form energy functional]
\label{thm:closed_form}
Let $\mu(x) = x/(x+1)$ and define $F(y)$ by the AQUAL
identity~\eqref{eq:aqual_identity} with boundary condition
$F(0) = 0$.  Then
\begin{equation}
  \boxed{F(y) = y - 2\sqrt{y} + 2\ln(1 + \sqrt{y}).}
  \label{eq:F_closed_form}
\end{equation}
\end{theorem}

\begin{proof}
The AQUAL identity~\eqref{eq:aqual_identity} together with
$F(0) = 0$ gives
\begin{equation}
  F(y) = \int_0^y \mu(\sqrt{t})\, dt
       = \int_0^y \frac{\sqrt{t}}{\sqrt{t} + 1}\, dt.
  \label{eq:integral_setup}
\end{equation}
Substitute $u = \sqrt{t}$, so that $t = u^2$ and $dt = 2u\, du$.
The limits transform as $t: 0 \to y$ becomes $u: 0 \to \sqrt{y}$.
Equation~\eqref{eq:integral_setup} becomes
\begin{equation}
  F(y) = \int_0^{\sqrt{y}} \frac{u}{u+1}\,(2u)\, du
       = 2\int_0^{\sqrt{y}} \frac{u^2}{u+1}\, du.
  \label{eq:after_sub}
\end{equation}
Apply the algebraic identity
\begin{equation}
  \frac{u^2}{u+1} = \frac{(u^2 - 1) + 1}{u+1}
                  = (u - 1) + \frac{1}{u+1},
  \label{eq:long_division}
\end{equation}
which follows from $u^2 - 1 = (u-1)(u+1)$.  Substituting
into~\eqref{eq:after_sub}:
\begin{equation}
  F(y) = 2\int_0^{\sqrt{y}}
         \left[(u - 1) + \frac{1}{u+1}\right] du.
\end{equation}
Integrate term by term using the elementary antiderivatives
$\int u\, du = u^2/2$, $\int du = u$, and
$\int (u+1)^{-1}\, du = \ln(u+1)$:
\begin{equation}
  F(y) = 2\left[\frac{u^2}{2} - u + \ln(u+1)\right]_0^{\sqrt{y}}.
\end{equation}
Evaluate at the upper limit $u = \sqrt{y}$:
\begin{equation}
  \frac{(\sqrt{y})^2}{2} - \sqrt{y} + \ln(\sqrt{y} + 1)
  = \frac{y}{2} - \sqrt{y} + \ln(1 + \sqrt{y}),
\end{equation}
and at the lower limit $u = 0$ each term vanishes
($0 - 0 + \ln 1 = 0$).  Therefore
\begin{equation}
  F(y) = 2 \cdot \left[\frac{y}{2} - \sqrt{y}
       + \ln(1 + \sqrt{y})\right]
       = y - 2\sqrt{y} + 2\ln(1 + \sqrt{y}).
\end{equation}
\hfill$\square$
\end{proof}

% ==================================================================

\section*{Verification}

\begin{lemma}[Recovery of $\mu$ by differentiation]
\label{lem:recovery}
Differentiating the closed form~\eqref{eq:F_closed_form} with
respect to $y$ recovers the simple interpolating function:
$dF/dy = \mu(\sqrt{y})$.
\end{lemma}

\begin{proof}
To keep all intermediate expressions regular at $y = 0$, work
in the natural variable $u = \sqrt{y}$ used in the proof of
Theorem~\ref{thm:closed_form}, with $y = u^2$ and $dy = 2u\, du$.
In $u$-coordinates the closed form~\eqref{eq:F_closed_form} reads
\begin{equation}
  F = u^2 - 2u + 2\ln(1 + u).
\end{equation}
Differentiating with respect to $u$ gives a sum of three regular
terms:
\begin{equation}
  \frac{dF}{du} = 2u - 2 + \frac{2}{1 + u}.
\end{equation}
Combining over the common denominator $(1+u)$:
\begin{equation}
  \frac{dF}{du}
  = \frac{2u(1+u) - 2(1+u) + 2}{1+u}
  = \frac{2u + 2u^2 - 2 - 2u + 2}{1+u}
  = \frac{2u^2}{1+u}.
\end{equation}
For $u > 0$, apply the chain rule
$dF/dy = (dF/du)/(dy/du) = (dF/du)/(2u)$:
\begin{equation}
  \frac{dF}{dy}
  = \frac{1}{2u} \cdot \frac{2u^2}{1+u}
  = \frac{u}{1+u}
  = \frac{\sqrt{y}}{1 + \sqrt{y}}
  = \mu(\sqrt{y}).
  \label{eq:recovery}
\end{equation}
The recovered $\mu$ matches the input~\eqref{eq:simple_mu}.  The
integration of Theorem~\ref{thm:closed_form} and the
differentiation here are mutual inverses on $u > 0$, as required
by the fundamental theorem of calculus.  Since the reduced
expression $u/(1+u)$ is regular at $u = 0$, it follows that
$\mu(\sqrt{y}) = \sqrt{y}/(1 + \sqrt{y})$ on all $y \geq 0$,
with $\mu(0) = 0$ recovered by direct substitution.
\hfill$\square$
\end{proof}

% ==================================================================

\section*{Corollaries}

\begin{corollary}[Newtonian limit]
\label{cor:newtonian}
In the high-acceleration regime ($y \gg 1$), $F(y) \to y$,
the corresponding response is $\mu \to 1$, and the modified
Poisson equation~\eqref{eq:modified_poisson} reduces to the
Newtonian $\nabla^2\Phi = 4\pi G\rho$.
\end{corollary}

\begin{proof}
For $y \gg 1$, $\sqrt{y} \gg 1$ and the subleading terms
$-2\sqrt{y} + 2\ln(1 + \sqrt{y})$ grow as $O(\sqrt{y})$, which
is negligible relative to $y$.  Hence $F(y) \sim y$ and
$dF/dy \to 1 = \mu(\infty)$.
\hfill$\square$
\end{proof}

\begin{corollary}[Deep-MOND limit]
\label{cor:deep_mond}
In the low-acceleration regime ($y \ll 1$), $F(y) \sim
\tfrac{2}{3}\, y^{3/2}$, which produces the asymptotic relation
$g_t \to \sqrt{g_s a_0}$ and the baryonic Tully-Fisher relation
$M \propto v^4$.
\end{corollary}

\begin{proof}
Expand $\ln(1 + \sqrt{y})$ in its Taylor series about
$\sqrt{y} = 0$:
\begin{equation}
  \ln(1 + \sqrt{y}) = \sqrt{y} - \frac{y}{2}
                    + \frac{y^{3/2}}{3} - \frac{y^2}{4}
                    + O(y^{5/2}).
  \label{eq:taylor}
\end{equation}
Substituting into~\eqref{eq:F_closed_form}:
\begin{equation}
  F(y) = y - 2\sqrt{y} + 2\sqrt{y} - y
       + \frac{2 y^{3/2}}{3} - \frac{y^2}{2} + O(y^{5/2})
       = \frac{2}{3}\, y^{3/2} + O(y^2).
  \label{eq:deep_mond}
\end{equation}
The leading behaviour $F(y) \sim (2/3)\, y^{3/2}$ is the
deep-MOND form~\cite{BekensteinMilgrom1984}.
\hfill$\square$
\end{proof}

\begin{corollary}[Smooth crossover]
\label{cor:crossover}
The transition between the Newtonian and deep-MOND regimes is
smooth: $F$ and all its derivatives are well-defined for
$y > 0$, in particular across the crossover scale $y \sim 1$.
No patching of two asymptotic forms is required.
\end{corollary}

\begin{proof}
The logarithmic term $2\ln(1 + \sqrt{y})$ is finite and
continuous for all $y \geq 0$.  Hence $F(y)$ is $C^\infty$ on
$(0, \infty)$ and continuous at $y = 0$.
\hfill$\square$
\end{proof}

% ==================================================================

\section*{Remarks}

\begin{remark}[Three-term structure]
\label{rem:three_terms}
The three terms of~\eqref{eq:F_closed_form} carry distinct
asymptotic roles.  At high $y$, the polynomial $y$ dominates and
recovers the Newtonian limit.  At low $y$, the polynomial,
square-root, and logarithmic terms enter into the Taylor
cancellation shown in Corollary~\ref{cor:deep_mond}: the linear
and square-root orders cancel among all three terms, leaving the
deep-MOND $y^{3/2}$ scaling as the surviving leading term.  The
coefficients $1$, $2$, $2$ and the appearance of exactly three
terms are not adjusted; they are forced by the integration of the
input $\mu(x) = x/(x+1)$.
\end{remark}

% ==================================================================

\begin{thebibliography}{9}

\bibitem{Milgrom1983}
M.~Milgrom,
\textit{A modification of the Newtonian dynamics as a possible
alternative to the hidden mass hypothesis},
Astrophys.\ J.\ \textbf{270} (1983) 365.

\bibitem{BekensteinMilgrom1984}
J.~Bekenstein and M.~Milgrom,
\textit{Does the missing mass problem signal the breakdown of
Newtonian gravity?},
Astrophys.\ J.\ \textbf{286} (1984) 7.

\bibitem{FamaeyMcGaugh2012}
B.~Famaey and S.~S.~McGaugh,
\textit{Modified Newtonian Dynamics (MOND): Observational
Phenomenology and Relativistic Extensions},
Living Rev.\ Relativ.\ \textbf{15} (2012) 10.

\end{thebibliography}

\end{document}
